% Part: second-order-logic
% Chapter: metatheory
% Section: second-order-arithmetic

\documentclass[../../../include/open-logic-section]{subfiles}

\begin{document}

\olfileid{sol}{met}{spa}

\olsection{二阶算术}

回想一下，皮亚诺算术理论 $\Th{PA}$ 包含 $\Th{Q}$ 的八条公理，
\begin{align*}
& \lforall[x][\eq/[x'][\Obj 0]]\\
& \lforall[x][\lforall[y][(\eq[x'][y'] \lif \eq[x][y])]]\\
& \lforall[x][(\eq[x][\Obj 0] \lor \lexists[y][\eq[x][y']])]\\ 
& \lforall[x][\eq[(x + \Obj 0)][x]]\\
& \lforall[x][\lforall[y][\eq[(x + y')][(x + y)']]]\\
& \lforall[x][\eq[(x \times \Obj 0)][\Obj 0]]\\
& \lforall[x][\lforall[y][\eq[(x \times y')][((x \times y) + x)]]]\\
& \lforall[x][\lforall[y][(x < y \liff \lexists[z][\eq[(z' + x)][y]])]]\\
\intertext{以及所有如下形式的语句：}
& (!A(\Obj 0) \land \lforall[x][(!A(x) \lif !A(x'))]) \lif \lforall[x][!A(x)].
\end{align*}
后者是一个“模式”，即一种能生成算术语言中无穷多语句的模板，每个公式~$!A(x)$ 对应其中一句。
我们称这个模式为（一阶的）\emph{归纳公理模式}。 
在\emph{二阶} 皮亚诺 算术~$\Th{PA^2}$ 中，归纳可以表述为单个语句。$\Th{PA^2}$
由上述前八条公理加上（二阶的）\emph{归纳公理}组成：
\[
\lforall[X][((X(\Obj 0) \land \lforall[x][(X(x) \lif X(x'))]) \lif \lforall[x][X(x)])].
\]
它说的是：如果论域的一个子集~$X$ 包含~$\Assign{\Obj{0}}{M}$，
并且对任意~$x \in \Domain{M}$，若 $X$ 包含 $x$ 则也包含其后继~$\Assign{\prime}{M}(x)$（即 $X$ “在后继下封闭”），
那么 $X$ 就包含论域中的所有元素（即 $X = \Domain{M}$）。

该归纳公理保证了任何满足它的结构，其论域~$\Domain{M}$ 中只包含公理要求存在的那些元素，
即对所有 $n \in \Nat$，~$\num{n}$ 的值。$\Th{PA^2}$ 的模型不包含非标准数。

\begin{thm}
\ollabel{thm:sol-pa-standard}
如果 $\Sat{M}{\Th{PA^2}}$，则 $\Domain{M} =
\Setabs{\Value{\num{n}}{M}}{n \in \Nat}$。
\end{thm}

\begin{proof}
令 $N = \Setabs{\Value{\num{n}}{M}}{n \in \Nat}$，并假设
$\Sat{M}{\Th{PA^2}}$。显然，对任意 $n \in \Nat$，
$\Value{\num{n}}{M} \in \Domain{M}$，因此 $N \subseteq \Domain{M}$。

现在证明反方向的包含关系。考虑变元指派 $s$，满足 $s(X) = N$。
根据假设，\begin{align*}
\Sat{M}{\lforall[X][((X(\Obj 0) \land \lforall[x][(X(x)
      \lif X(x'))]) \lif \lforall[x][X(x)])]}, & ，\text{从而} \\
\Sat{M}{(X(\Obj 0) \land \lforall[x][(X(x) \lif X(x'))]) \lif
  \lforall[x][X(x)]}[s]. & 
\end{align*}。
考虑这个条件句的前件。首先，$\Value{\Obj{0}}{M} \in
N$，因此 $\Sat{M}{X(\Obj{0})}[s]$。
第二个合取支，即 $\lforall[x][(X(x) \lif X(x'))]$，同样被满足。
因为若 $x \in N$，根据 $N$ 的定义，对某个 $n$ 有 $x = \Value{\num{n}}{M}$。
由此可得 $\Assign{\prime}{M}(x) = \Value{\num{n+1}}{M} \in N$。所以，
$\Assign{\prime}{M}(x) \in N$。

这样我们就得到 $\Sat{M}{X(\Obj 0) \land \lforall[x][(X(x) \lif X(x'))]}[s]$。
因此，$\Sat{M}{\lforall[x][X(x)]}[s]$。而这意味着，对于
每一个 $x \in \Domain{M}$，都有 $x \in s(X) = N$。所以，$\Domain{M}
\subseteq N$。
\end{proof}

\begin{cor}
\ollabel{cor:sol-pa-categorical}
$\Th{PA^2}$ 的任意两个模型都是同构的。
\end{cor}

\begin{proof}
根据 \olref{thm:sol-pa-standard}，$\Th{PA^2}$ 的任何模型的论域都被 $\Value{\num{n}}{M}$ 穷尽。
任何这样的模型也都是 $\Th{Q}$ 的模型。
再由 \olref[mod][mar][stm]{prop:thq-standard}，任何这样的模型都是标准的，即同构于 $\Struct{N}$。
\end{proof}

前面我们将 $\Th{PA^2}$ 定义为包含前八条算术公理加上二阶归纳公理的理论。事实上，借助
二阶逻辑的表达力，二阶 皮亚诺 算术只需要前两条算术公理加上归纳公理。

\begin{prop}\ollabel{prop:sol-pa-definable}
令 $\Th{PA^{2\dagger}}$ 为包含前两条算术公理（后继公理）和二阶归纳公理的二阶理论。
则 $\le$、$+$ 和 $\times$ 在 $\Th{PA^{2\dagger}}$ 中是可定义的。
\end{prop}

\begin{proof}
要证明 $\le$ 可定义，我们需要找到一个公式 $!A_\le(x, y)$，
使得 $\Sat{N}{!A_\le(\num n, \num m)}$ 当且仅当 $n \le m$。考虑公式
\[
!B(x, Y) \ident Y(x) \land \lforall[y][(Y(y) \lif Y(y'))]
\]
显然，一个集合 $Y \subseteq \Nat$ 满足 $!B(\num n, Y)$，当且仅当
$\Setabs{m}{n \le m} \subseteq Y$，因此我们可以取 $!A_\le(x, y) \ident
\lforall[Y][(!B(x, Y) \lif Y(y))]$。

要说明加法可定义，注意到 $k+l=m$ 当且仅当存在一个函数 $u$，使得 $u(0)=k$，$u(n')=u(n)'$ 对所有 $n$ 成立，并且 $m = u(l)$。 
我们可以利用这个等价性，用以下公式在 $\Th{PA^{2\dagger}}$ 中定义加法：
\[
!A_+(x,y,z) \ident \lexists[u][(u(\Obj 0)=x \land \lforall[w][u(x')=u
 (x)'] \land u(y)=z)]
\] 
不难看出，$\Sat{N}{!A_+(\num k, \num l, \num m)}$ 
当且仅当 $k+l=m$。
\end{proof}

\begin{prob}
完成 \olref[sol][met][spa]{prop:sol-pa-definable} 的证明。
\end{prob}

\end{document}